\chapter*{Liste des abréviations}
\addcontentsline{toc}{chapter}{Liste des abréviations}
\markboth{Liste des abréviations}{}

\begin{table}[H]
    \centering
    \small
    \renewcommand{\arraystretch}{0.92}
    \caption{Liste des abréviations utilisées dans le rapport}
    \label{tab:abreviations}
    \begin{tabular}{@{}>{\raggedright\arraybackslash}p{3cm} >{\raggedright\arraybackslash}p{11cm}@{}}
        \toprule
        \textbf{Abréviation} & \textbf{Signification} \\
        \midrule
        API & Application Programming Interface (interface de programmation applicative) \\
        ARP & Address Resolution Protocol (protocole de résolution d'adresse) \\
        CIDR & Classless Inter-Domain Routing (notation de plage d'adresses réseau) \\
        CLI & Command-Line Interface (interface en ligne de commande) \\
        EMSI & École Marocaine des Sciences de l'Ingénieur \\
        HTTP & HyperText Transfer Protocol \\
        IP & Internet Protocol \\
        IT & Information Technology (technologies de l'information) \\
        JSON & JavaScript Object Notation \\
        JWT & JSON Web Token \\
        MAC & Media Access Control (adresse physique d'une interface réseau) \\
        MCD & Modèle Conceptuel de Données \\
        MVP & Minimum Viable Product (produit minimum viable) \\
        OSI & Open Systems Interconnection (modèle de référence réseau) \\
        OUI & Organizationally Unique Identifier (préfixe MAC identifiant le fabricant) \\
        REST & Representational State Transfer (style d'architecture d'API web) \\
        SPA & Single Page Application \\
        SQL & Structured Query Language \\
        TCP & Transmission Control Protocol \\
        UML & Unified Modeling Language (langage de modélisation unifié) \\
        \bottomrule
    \end{tabular}
\end{table}
